\section{Data}

\subsection{Data Sources}

We focused our attention on a few public databases having information on basic classification of diseases (MeSH),
 drug chemical descriptors and targets (ChEMBL), disease targets (Open Targets) as well as information on
  drug--disease indication studies (ClinicalTrials).
 Table~\ref{tab:data_sources} summarizes information on these data sources. Although initially we were using the web
 scraping to get information on disease classification, we found out that this way of getting medical 
 information is rather deprecated, so we queried the corresponding databases instead. Data concerning drugs and diseases
were combined using the cross references furnished by the same databases (ChEMBL and Open Targets). We also used a UniProt identifier mapper
that allowed us to unify synonymous target IDs as well as enrich them with corresponding 
molecular pathway identifiers from Reactome database.

\begin{table}[H]
    \centering
    \caption{Description of data sources.}
    \label{tab:data_sources}
    \rowcolors{2}{gray!15}{white} % alternate row colors
    \begin{tabular}{p{3cm} p{5cm} p{6cm}}
        \toprule
        \textbf{Source} & \textbf{Access Method} & \textbf{Data Retrieved} \\
        \midrule
        MeSH RDF & SPARQL endpoint & Disease IDs and ontology \\
        ChEMBL & REST API (Python client with query-style filters) & Drug chemical properties, targets and indications\\
        UniProt & REST API (Python client) & Map target and pathway IDs\\
        Open Targets & GraphQL API  & Disease targets\\
        ClinicalTrials.gov & REST API & Trial results and metadata\\
        \bottomrule
    \end{tabular}
\end{table}

\subsection{Target variable definition}

The prediction target is a success--fail variable indicating whether a drug shows an effect against a disease;
this constitutes a binary label for each drug--disease indication. Drugs that passed
the fourth (final) clinical trial phase for a given disease are labeled as successful.
For other cases, the outcome is less clear: if the last phase wasn't (yet) passed, we can understand the reason
by addressing the corresponding clinical trials. There are several possibilities: the combination was never tried (most common), was tried but ineffective, was effective only in earlier phases,
 or has no published results due to study termination or other reasons.
There exists some manually and automatically annotated clinical trial databases \cite{gao2024automatic} \cite{chen2025trialbench},
 but they are either limited in the time span (e.g. last 10 years), or their "non-success" label simultaneously encompasses
trials failed due to drug inefficiency and other reasons. Therefore, we decided to automatically 
label drug--disease indications that have been tried at least once using rules detailed in Algorithm~\ref{alg:labeling} 
(the complete algorithm can be found in the source code in the data preparation section).

\begin{figure}[H]
    \centering
   \includegraphics[width=0.7\textwidth]{figures/03-label-stability_in_time.png}
    \caption{Counts and proportions of drug--disease indications by success label over time.}
    \label{fig:03-label-stability_in_time}
\end{figure}

Interestingly, most undecided indications and clinical trials
without published results occur roughly before 2000 and after 2010, as shown in Figures~\ref{fig:03-label-stability_in_time}~and~\ref{fig:03-label-has_results_in_time}. One can speculate
that it is linked to the fact that ClinicalTrials.gov was launched in 2000 and it 
became mandatory to publish there all studies (but not full results) since 2007. As can be seen in Figure~\ref{fig:03-label-distribution_missing_date}, 
among indications with no first trial start date available, most are labeled successful, 
which mostly corresponds to indications established long before the launch of ClinicalTrials.gov. 
It's important to note that some of those choices are arbitrary, which likely constitutes 
the most significant limitation.

\subsection{Feature Engineering and Description}

Detailed description of extracted features can be found in Table~\ref{tab:feature_description}.
Most drug features describe chemical and structural properties. Among engineered features are Morgan molecular fingerprints with 2 atoms radius 
obtained with RDKit module from drug's SMILES (flat string representation of the structure) and yielding a 2048 dimensional bit vector. Each
bit of this vector encodes the presence or absence of a particular substructure. To reduce the dimensionality, 
first 100 dimensions were taken and labeled as \verb|FP_#|.

Disease features consist of disease ontology description, one of them being ontological category
in addition to "immune system diseases", labeled as \verb|cat_#|, for example some of immune diseases are as well 
neoplastic (e.g. leukemia).

For both diseases and drugs, we obtained a list of molecular targets, focusing on proteins. 
Since targets themselves are highly heterogenous and deeply interconnected
within molecular pathways, we converted each drug's or disease's target list into a corresponding list of pathways. 
We then applied TF-IDF vectorizations separately to the pathway lists for drugs and for diseases. 
This approach downweights common housekeeping pathways while highlighting more informative ones. 
We then performed PCA and extracted first 50 principal components for drugs and 50 for diseases, labeled as \verb|drug_path_#| and \verb|disease_path_#|. 
 This information captures drug-drug and disease-disease mechanistic similarity based on shared pathways.

 To capture the drug--disease pathway similarity, we computed a Jaccard similarity index based on corresponding pathway 
 ID lists. We also considered including drug--disease name similarity based on PubMedBERT embeddings, which reflect  how closely a drug and a disease are associated in the scientific literature.
 However, we later realized that including this feature in model training could introduce a data leakage risk, 
 since PubMedBERT is trained on recent literature.


\subsection{Exploratory Data Analysis}

\subsubsection{Missing Values}

Among collected data and selected features, only a specific group of drug features contains missing 
values (Figure~\ref{fig:03-na-drugs_with_bioth}).
To investigate the cause, we created an indicator variable for missingness and visualized patterns with a correlation heatmap (Figure~\ref{fig:03-corr_matrix-na}). 
Indeed, this group of chemical descriptors (such as number of aromatic rings or lipophilicity) tends to be absent together
 and is associated with biotherapeutics marker. Biotherapeutics are drugs derived from living 
 organisms rather than being chemically synthesized. Most of them are antibodies or other proteins - 
 large biopolymers for which these chemical descriptors are not relevant. 

\begin{figure}[H]
    \centering
    \begin{minipage}[t]{0.48\textwidth}
        \centering
        \includegraphics[width=\textwidth]{figures/03-na-drugs_with_bioth.png}
        \caption{Missing values heatmap in the drugs dataset.}
        \label{fig:03-na-drugs_with_bioth}
    \end{minipage}
    \hfill
    \begin{minipage}[t]{0.48\textwidth}
        \centering
        \includegraphics[width=\textwidth]{figures/03-corr_matrix-na.png}
        \caption{Missingness analysis of drug descriptors with correlation heatmap.}
        \label{fig:03-corr_matrix-na}
    \end{minipage}
\end{figure}

Inspection of drugs with missing values confirms this (e.g. insulin appears among such drugs). Therefore, we restricted
the analysis to small-molecule drugs,  since chemical descriptors are not applicable to biopolymers and imputing 
 them would be inappropriate. After removing these cases, the data was ready for merging, 
 resulting in a dataset containing complete information for each drug--disease indication. Table~\ref{tab:03-dataset_summary} provides 
 a summary of data dimensions and label counts.

\begin{table}[H]
\centering
\caption{Summary of data dimensions and label counts.}
\label{tab:03-dataset_summary}
\rowcolors{2}{gray!15}{white} 
\input{tables/03-dataset_summary.tex}
\end{table}


\subsubsection{Univariate and Bivariate Analysis}

Inspection of feature distributions across label values for variables unrelated to molecular pathways or fingerprints 
(Figure~\ref{fig:03-feature_dist_by_success}) shows mostly bell-shaped distributions, sometimes multimodal or slightly skewed. Log transformation was
therefore considered unnecessary. Some disease ontology variables and drug descriptors exhibit outliers,
 which likely reflect normal biological or chemical variability and were further examined with dimension reduction techniques.
Overall, feature distributions are similar between success and failure groups, as expected given the limited discriminative power of
 basic descriptors alone. However, engineered name and path similarity features tend to be higher for 
successful indications, suggesting potential predictive value when combined with
other variables. Name similarity showed similar trend but was excluded from further analysis due to potential information leakage.

On a correlation heatmap of the full merged dataset (Figure~\ref{fig:03-corr_matrix-all_features}), reveals, as expected, clusters of strongly
 or moderately associated variables. Notably between chemical descriptors, molecular fingerprints, and drug pathway embeddings,
  as well as between disease ontology types and disease pathway embeddings.

 Focusing on more interpretable features only ( Figure~\ref{fig:03-corr_matrix-short}), allows straightforward 
 interpretation of some associations. For instance, the number of heavy atoms is associated with molecular weight. Number hydrogen bond acceptors and donors (HBA and HBD) 
correlate positively with with  polar surface area (PSA) and negatively with octanol-water partition coefficient (AlogP),
 reflecting increased hydrophilicity. Number of Lipinski rule-of-five violations is consistently associated with these descriptors.
  Interestingly, the number of aromatic rings correlates negatively with the natural product-likeness score (NP-likeness score), suggesting that
   purely synthetic drugs tend to contain more aromatic rings. Regarding disease ontology, immune diseases labelled as hemic or lymphatic tend 
to be classified as neoplastic, which is consistent with the fact that most immune-related cancers are leukemias or lymphomas.
 Another example is the positive association between urogenital diseases and infections.

\begin{figure}[H]
    \centering
   \includegraphics[width=0.9\textwidth]{figures/03-corr_matrix-short.png}
    \caption{Focused view of the correlation heatmap on interpretable features.}
    \label{fig:03-corr_matrix-short}
\end{figure}

\subsubsection{Dimension Reduction and Outlier Analysis}

PCA performed on centered and scaled merged data (Figure~\ref{fig:03-pca-scatter_plot}) shows that first two principal components explain only about 10 \%
of total variance, with variability spread across many components, as indicated by the scree plot (Figure~\ref{fig:03-pca-scree_plot}).
 This suggests the absence of of dominant linear axes of variation, consistent with an intrinsically
high-dimensional feature space and likely non-linear associations in such a complex system.
Furthermore, many features were explicitly constructed to be uncorrelated, such as principal components derived from molecular fingerprints and pathway embeddings.
Nevertheless, inspection of the projection onto the first two components reveals 
that the main source of variation is driven by disease heterogeneity. In particular, most influential variables in the correlation circle (Figure~\ref{fig:03-pca-corr_circle}) 
correspond to disease category indicators and disease pathway embeddings, distinguishing, for example, neoplastic blood disorders from respiratory tract diseases. 


\begin{figure}[H]
    \centering
    \begin{minipage}[t]{0.48\textwidth}
        \centering
        \includegraphics[width=\textwidth]{figures/03-pca-scatter_plot.png}
        \caption{PCA projection of samples onto the first two principal components.}
        \label{fig:03-pca-scatter_plot}
    \end{minipage}
    \hfill
    \begin{minipage}[t]{0.48\textwidth}
        \centering
        \includegraphics[width=\textwidth]{figures/03-pca-corr_circle.png}
        \caption{PCA correlation circle showing feature contributions to the first two principal components. Only arrows with length above 0.5 are shown.}
        \label{fig:03-pca-corr_circle}
    \end{minipage}
\end{figure}

The second principal component is primarily associated with drug-related variability. In particular,
 a small group of potential outliers appears at high values of both first and second principal components.
Closer inspection reveals that these indications involve three drugs - sirolimus, everolimus and ridaforolimus,
which are antineoplastic and immunosuppressive drugs sharing a common structural core - large 
macrolide ring (Figure~\ref{fig:03-rdkit_structure-drug_outliers}). We therefore interpret these observations 
as a part of normal chemical variability rather than errors or inadequate filtering and chose not to exclude them.


Although PCA is an unsupervised method, a posteriori labeling reveals a slight separation between success and
 failure cases along the diagonal defined by the first two principal components. Notably, the Quantitative Estimate of Druglikeness (QED)
  contributes to both components and is aligned with successful indications. This suggests that the dominant directions of
   variability capture, at least partially, differences related to indication success. This pattern may also reflect differences in success rates 
   across disease types; for instance, clinical studies in neoplastic blood diseases may exhibit lower apparent success rates due to higher study frequency 
   or greater clinical difficulty.

Non-linear dimension reduction methods such as UMAP and t-SNE (Figures~\ref{fig:03-umap}~and~\ref{fig:03-tsne}) reveal
 a large number of small clusters.These clusters appear to correspond to the indications involving identical or closely related drugs or diseases.
 In contrast, Isomap (Figure~\ref{fig:03-isomap}) does not reveal any easily interpretable structure in the data.

\subsubsection{Supervised Dimension Reduction}
Supervised dimension reduction using LDA and PLS-DA (Figures~\ref{fig:03-lda_projection}~and~\ref{fig:03-plsda}) 
shows that, although overlapping, success and failure labels are relatively well separated along the LDA axis 
and in projection onto the first two PLS-DA components.

\begin{figure}[H]
    \centering
    \begin{minipage}[t]{0.48\textwidth}
        \centering
        \includegraphics[width=\textwidth]{figures/03-lda_projection.png}
        \caption{Linear Discriminant Analysis (LDA) projection highlighting class separation.}
        \label{fig:03-lda_projection}
    \end{minipage}
    \hfill
    \begin{minipage}[t]{0.48\textwidth}
        \centering
        \includegraphics[width=\textwidth]{figures/03-plsda.png}
        \caption{Partial Least Squares Discriminant Analysis (PLS-DA) projection of samples highlighting class separation.}
        \label{fig:03-plsda}
    \end{minipage}
\end{figure}

\subsection{Graph Representation}


Although our binary classification approach for drug--disease indications is valid, 
it does not fully capture the true structure of the problem. The data are inherently structured as a 
bipartite drug--disease graph, with connections mediated by associated molecular pathways. Most advanced
 drug discovery research leverages molecular interaction networks and Graph Neural Networks directly, rather than relying 
 on manually engineered features \cite{keramida2025ai}. To explore this, we organized the same data as 
 a \verb|HeteroData| graph using the PyTorch Geometric module (Listing~\ref{lst:heterodata}). 
 Building a robust graph-based model would likely require richer information about molecular and mechanistic connections, 
 as well as a larger number of observations.


